\section{Knowledge organization and lifecycle management}
\label{sec:organization}

Acquisition decides what enters the system; organization decides what the
system is. Here organization is a governed process rather than a folder
convention or an emergent property of embedding similarity, with three
invariants: content decides where a document belongs, the disk, the tree
registry and the library metadata must agree about every document, and every
operation that removes or relocates knowledge is preceded by a survey and
followed by a verification. This chapter describes their implementation and
where the three stores can still diverge.

\subsection{Content-routed placement}
\label{sec:org-placement}

Placement reads the parsed document, not its filename or its folder of
origin. The pipeline selects the library whose stated scope covers the
content, or proposes a new library when no existing scope does, writes the
document, derives a description and a tag set, and runs the same final check
used by acquisition before the document counts as placed. Two rules keep the
catalog usable as a retrieval instrument: an existing library is preferred
over a new one; and a document is placed once, because re-routing is an
explicit lifecycle operation rather than a side effect of a later ingest.

The decision itself is made over library scope rather than folder names. The
ingest agent is given the candidate libraries with their scope descriptions and
the parsed body of the document, and selects the library whose scope covers the
content. That assignment is recorded in the catalog and the document is written
to that library. The decision is per document, not per part: a split document's
parts inherit the assignment made for their source. When no scope fits, the
agent defines a new library and states its scope before writing into it. The one
ingest run we instrumented carried its assignments in a classification file
keyed by paper identifier and built from their content together with the library
descriptions; an unassigned paper raised an error rather than a guess. The
prototype fixes neither a model nor a similarity threshold for the choice, so it
remains operator-visible at review time.

The library tree supports one level of content grouping in which child
libraries are addressable through a parent. The implementation carries a
caveat we report rather than hide: the vector segments of child documents are
stored in the parent library's collection, so a search that targets a child
library identifier returns no candidates. The shipped operational guidance is
to search the parent collection and to use the child libraries for navigation
and scoping rather than as search endpoints. This is a limitation of the
storage mapping, not of the routing rule.

\subsection{Catalog and metadata registries}
\label{sec:org-registries}

Table~\ref{tab:datamodel} lists the fields of the two registries that
describe the same collection: the global tree (\texttt{.tree-fs.json}) is the
authoritative structure, and each library carries a
\texttt{.knowledge-base.yml} indexing its Markdown documents and their image
directory.

Agents read the catalog through a progressive-disclosure interface: the
listing returns names, descriptions and counts, document listings return
metadata without bodies, and only an explicit read returns content,
attributed to a specific document. A caller can therefore narrow a search to
a handful of candidates without loading a body, which makes the catalog lane
viable at 5--6 turns and 54--91~s (Table~\ref{tab:lanes}).

Correctness of the catalog is three-source consistency: for every document,
the file on disk, its entry in the tree registry, and its record in the
library metadata must agree on existence, path, size and index state. The
platform verifies this agreement rather than assuming it: in the acceptance
campaign a renamed document kept its old vector segments, so the recorded
chunk count read 18 where 9 chunks existed, and the document matched its own
former path at 0.9738 similarity, high enough to be reported as a
near-duplicate pair. Only an explicit force re-index repaired it
(18~$\rightarrow$~9). Metadata and index consistency are therefore two
different obligations: the first is satisfied by a single write, the second
is not, because the vector store is updated by a separate call.

\subsection{Tag taxonomy and duplicate detection}
\label{sec:org-tags}
\label{sec:org-dedup}

Tags are the coarse, human-legible axis of the collection, and retrieval by
tag is a first-class operation. A tag registry holds the vocabulary; each
document's tags are normalized, filtered against a blacklist of structural
strings (section headings, generic words such as ``introduction''), and
re-checked against the document body so that a tag must occur in the content
it labels. Governance is a maintenance loop rather than a schema: no taxonomy
is imposed on authors and the vocabulary grows with the corpus, so the audit
reports tags with zero referencing documents, accumulated from test documents
and chapter headings, and removes them only when instructed. In the audited
campaign 165 documents passed the tag gate with no failures, and the live
registry held 316 tags. We do not claim a controlled vocabulary; we claim
that the vocabulary is inspectable and that orphaned entries are detectable.

Deduplication is a collection-level question as well: the same content can
arrive through two ingest paths or into two libraries. The procedure is
survey-then-act --- enumerate candidate duplicate groups at a similarity
threshold (0.90, with a manual fallback channel at 0.85), inspect each group,
delete nothing automatically --- and the audited run illustrates why the
survey step matters more than the threshold. On the 322-document corpus the
detector reported 35 candidate groups with similarities between 0.90 and
0.975, every one of them a pair of adjacent parts of the same long paper, and
no deletion was performed; a policy that deleted on detection would have
destroyed 35 legitimate document parts. The same run confirmed zero true
duplicates. Cross-library duplication is detected but not merged
automatically, because merging two libraries carries provenance consequences
and stays with the operator.

\subsection{Library lifecycle operations}
\label{sec:org-lifecycle}

Table~\ref{tab:org-lifecycle} lists the operations the platform exposes, the
interface that performs them, the consistency obligation each one carries,
and the state observed for each.

\begin{table}[t]
\centering
\small
\caption{Library and document lifecycle operations, their consistency
obligation, and the state observed in the acceptance campaign. ``Audited''
means the three stores were compared after the operation.}
\label{tab:org-lifecycle}
\begin{tabular}{@{}p{0.22\textwidth}p{0.24\textwidth}p{0.20\textwidth}p{0.24\textwidth}@{}}
\toprule
Operation & Interface & Obligation & Observed \\
\midrule
Create library & library create & tree + metadata + directory & audited \\
Rename or move document & document move & path updated in all three stores & audited; stale segments until force re-index \\
Update description or tags & document update & metadata rewritten, body re-checked & audited (165/165) \\
Update content & document update & new hash, re-index required & re-index not implicit \\
Delete document & document delete & removal from all three stores & audited (1~$\rightarrow$~0) \\
Batch delete & batch delete & per-document verification & 1/1 in campaign \\
Move library & library move & subtree and metadata move together & audited (0~$\rightarrow$~1) \\
Merge libraries & library merge & provenance preserved & demonstrated, not controlled \\
Split library & library split & no part orphaned & analysed, deliberately not executed \\
Rebuild index & re-index (task id) & segment count matches content & 18~$\rightarrow$~9 after force \\
\bottomrule
\end{tabular}
\end{table}

The rule the table carries is that re-indexing is not implicit: creating,
moving or rewriting a document does not update the vector store, only an
explicit index operation does, and the audit treats a missing or stale index
as an incomplete operation rather than as a warning. The contribution here is
a contract under which disagreement between an approximate index and a
mutating document store becomes detectable, not a new maintenance algorithm
\citep{xu-spfresh-2023}. The evidence boundary is narrower: move, rename,
update, delete and re-index were exercised with real tool calls and audited
afterwards; merge and split are implemented and were demonstrated on live
libraries, but they have no controlled evaluation, and the sub-library
hierarchy was analysed and deliberately not executed pending an operator
decision. No concurrent organizational operation was tested.

\subsection{The document graph}
\label{sec:org-graph}

The platform maintains a property graph with document, library, tag, folder
and entity nodes connected by membership, tagging, relatedness, identity and
containment edges. Those edges derive from tag co-membership, vector
similarity above a threshold, and the verdicts produced by content
verification during retrieval. Adding one document to a library added four
edges in the audited run (4{,}675~$\rightarrow$~4{,}679), one per assigned
tag; a full build over the 100-paper corpus produced 658 nodes covering 333
documents and 316 tags, and the 165-document intermediate state held 335
nodes and 1{,}671 edges.

The graph keeps a deliberately restricted role. It is a navigation instrument
that expands a candidate set and surfaces documents sharing context across
libraries. It is never evidence: as a standalone retriever it reached Hit@5
of 0.5625 on the industrial corpus and 0.2 on the second corpus, well below
the lexical and dense channels (Table~\ref{tab:retrieval}). Graph-based
organization in this system is useful for discovery and weak as a ranker.

Two implementation gaps belong here. Graph construction is not yet guaranteed
by the ingestion path: a document ingested through one agent harness received
no graph node, so a library reported two of three documents. And batch graph
builds have been observed to leave a task record in the running state after
the work completed, so task status cannot be used as completion evidence; the
graph itself must be queried.

\subsection{Operational knowledge as a first-class collection}
\label{sec:org-experience}

Process documents describe what the organization intends; operational lessons
describe what actually happened. The platform stores the second kind as a
first-class collection with the same addressing, indexing and retrieval
interface as ordinary documents, but with two differences that reflect the
epistemic status of a lesson. First, an experience carries a credibility tier
that controls whether it may be shown (Table~\ref{tab:exp-tiers}); the
experience channel's similarity threshold is 0.55 against the document
channel's 0.35, and an unvetted or disputed entry is capped even when it
scores well. Second, entries are drafted before they are served: extraction
produces candidates from ingested content or from recurring question-answer
pairs, and automated distillation runs as a separate agent process on a
schedule, guarded by a circuit breaker (three consecutive failures suspend it
for 24 hours) and a heuristic fallback marked low-confidence. Nothing reaches
the served collection without an explicit approval step, the same fail-closed
posture the retrieval gate applies to citations.

\begin{table}[t]
\centering
\small
\caption{Experience credibility tiers and modifiers. Tiers combine retrieval
score, content-verification score, ratings and review count; modifiers cap
the tier rather than promote it.}
\label{tab:exp-tiers}
\begin{tabular}{@{}l>{\raggedright\arraybackslash}p{3.6cm}p{0.46\textwidth}@{}}
\toprule
Tier & Condition & Effect \\
\midrule
P0 strong & sim $\ge 0.65$ $\wedge$ content $\ge 6$ $\wedge$ rating $\ge 4$ $\wedge$ reviews $\ge 1$ & may be presented first \\
P1 reference & sim $\ge 0.45$ $\wedge$ content $\ge 4$ & presented with credibility noted \\
P2 weak & sim $\ge 0.35$ $\wedge$ content $\ge 3$ & suppressed unless expanded \\
discard & content verification fails or sim $< 0.35$ & never presented \\
\addlinespace
disputed & reviews $\ge 3$ $\wedge$ rating $< 2.0$ & capped at P2 \\
unvetted & 0 reviews $\wedge$ 0 applications & capped at P1 \\
auto-extracted & not human-reviewed & capped at P1 until applied twice and rated $\ge 4$ \\
\bottomrule
\end{tabular}
\end{table}

The evidence for this subsystem is deliberately case evidence. One extraction
campaign produced thirteen candidate drafts; one end-to-end scenario
sedimented a single experience, which a later global experience search
returned first with a local similarity of 0.711; and the approval path was
never exercised, because the draft pool was empty when it was tested. A
single scenario and thirteen candidates cannot support an effectiveness
claim, and we do not make one. What the campaign establishes is that
operational knowledge travels through the same pipeline as documents,
addressed, indexed, gated and auditable, rather than living in a wiki or a
prompt.

\subsection{Integrity verification as routine}
\label{sec:org-integrity}

Every property claimed in this chapter is checked by a verification pass
against the live collection, so that consistency, integrity and health are
measured rather than asserted. The pass is read-only by default; repairs
require an explicit instruction, and the report is produced first, because an
audit that repairs as it goes cannot report what it found.

The pass is not ceremonial. It is how the two index-integrity incidents of
this deployment were found, a silently lost collection and a stale index, and
how the reorganization campaign was certified after the fact: 322/322
metadata checks, 165/165 description checks, 33/33 sampled consistency
checks, and 15/15 retrieval regression probes with no degradation. In the
same campaign the audit surfaced the failures reported above: the rename-time
chunk duplication, the missing graph node, and a task-status record that
outlived its task. A collection whose health is a command is a collection
whose health can be argued about with evidence.

\subsection{What organization buys, and what it does not}
\label{sec:org-summary}

Organization buys candidate quality and auditability. Because the catalog is
curated and consistent, an agent can select a library and a document set
without loading bodies, which is the mechanism behind the catalog lane and
behind cross-library discovery in an environment with nineteen libraries.
Because every placement, tag, description and deletion left a record, the
collection can be re-verified after the fact; the numbers in this chapter are
audit outputs rather than survey estimates.

Three things organization does not do, and the paper does not claim. It does
not repair a bad parse: a document whose parse lost its tables cannot be
recovered by placing it well. It does not prevent divergence by construction:
the write path can leave the vector index behind, and only an audit and an
explicit re-index restore agreement. And it does not scale without limits:
the largest collection we report is 322 documents in five libraries, and the
merge and split operations are unevaluated. These are engineering boundaries.
